% ===== PRÉAMBULE CANONIQUE — à coller VERBATIM en tête de CHAQUE .tex =====
\documentclass[11pt,a4paper]{article}
\usepackage[utf8]{inputenc}\usepackage[T1]{fontenc}\usepackage{lmodern}
\usepackage[french]{babel}
\usepackage{amsmath,amssymb,mathtools}\usepackage{icomma}
\usepackage[version=4]{mhchem}          % \ce{...} : équations & formules
\usepackage{siunitx}\sisetup{locale=FR} % \qty{}{}, \si{}, \ang{}
\usepackage{chemfig}                    % \chemfig{...} : structures développées
\usepackage{xcolor}\usepackage{colortbl}
\usepackage{tikz}\usepackage{pgfplots}\pgfplotsset{compat=1.18}
\usepackage[most]{tcolorbox}\usepackage{enumitem}\usepackage{array,booktabs}
\usepackage[a4paper,margin=2.2cm]{geometry}
\usetikzlibrary{arrows.meta,calc,angles,quotes,positioning,patterns,backgrounds,shapes.geometric}
\definecolor{primary}{RGB}{222,49,99}\definecolor{primarydark}{RGB}{150,22,55}
\definecolor{defcol}{RGB}{0,114,178}\definecolor{methcol}{RGB}{52,92,148}
\definecolor{excol}{RGB}{0,135,90}\definecolor{attncol}{RGB}{200,40,55}
\tcbset{boxbase/.style={enhanced,breakable,boxrule=0.9pt,arc=2.5pt,
  left=3mm,right=3mm,top=2.3mm,bottom=2.3mm,fonttitle=\bfseries,coltitle=white}}
% --- boîtes TUTORIEL (noms EXACTS, ne pas renommer) ---
\newtcolorbox{cle}[1][]{boxbase,colback=defcol!6,colframe=defcol,
  title={\textsc{À retenir}\ifx\relax#1\relax\relax\else\textmd{ : \itshape #1}\fi}}
\newtcolorbox{methode}[1][]{boxbase,colback=methcol!7,colframe=methcol,sharp corners,
  title={\textsc{Méthode}\ifx\relax#1\relax\relax\else\textmd{ --- \itshape #1}\fi}}
\newtcolorbox{exemplebox}[1][]{boxbase,colback=excol!5,colframe=excol,
  title={\textsc{Exemple}\ifx\relax#1\relax\relax\else\textmd{ : \itshape #1}\fi}}
\newtcolorbox{piege}[1][]{boxbase,colback=attncol!6,colframe=attncol,
  title={\textsc{Piège}\ifx\relax#1\relax\relax\else\textmd{ : \itshape #1}\fi}}
\newtcolorbox{remarque}[1][]{boxbase,colback=black!4,colframe=black!45,
  title={\textsc{Remarque}\ifx\relax#1\relax\relax\else\textmd{ : \itshape #1}\fi}}
% --- boîtes EXERCICE / CORRIGÉ (noms EXACTS, auto-comptées) ---
\newtcolorbox[auto counter]{exo}[1][]{boxbase,colback=primary!4,colframe=primary,
  title={\textsc{Exercice}~\thetcbcounter\ifx\relax#1\relax\relax\else\textmd{ --- \itshape #1}\fi}}
\newtcolorbox[auto counter]{corrige}[1][]{boxbase,colback=excol!4,colframe=excol,
  title={\textsc{Corrigé}~\thetcbcounter\ifx\relax#1\relax\relax\else\textmd{ --- \itshape #1}\fi}}
\newcommand{\R}{\mathbb{R}}\newcommand{\N}{\mathbb{N}}\newcommand{\Z}{\mathbb{Z}}\newcommand{\ds}{\displaystyle}
% (un \usepackage{fancyhdr}+en-tête est possible ; il est ignoré côté web)
% =========================================================================
\begin{document}
\begin{center}\color{primarydark}\large\bfseries Thème 2 — Demi-équations et équilibrage \quad$\bullet$\quad Niveau Moyen\end{center}

\begin{exo}[le piège du « rapport 1 pour 1 »]
Équilibre chaque demi-équation en milieu acide. Attention au nombre d'atomes qui changent de n.o.
\begin{enumerate}[label=\alph*)]
  \item \ce{Cr2O7^{2-} <=> Cr^{3+}}
  \item \ce{MnO4^- <=> Mn^{2+}}
  \item \ce{SO4^{2-} <=> S} \quad (le soufre passe de $+\mathrm{VI}$ à $0$)
\end{enumerate}
\end{exo}

\begin{exo}[équation globale en milieu acide]
Le permanganate oxyde l'ion ferreux : couples \ce{MnO4^-}/\ce{Mn^{2+}} et \ce{Fe^{3+}}/\ce{Fe^{2+}}.
\begin{enumerate}[label=\alph*)]
  \item Écris les deux demi-équations équilibrées.
  \item Combine-les en équation globale.
  \item Donne la somme des coefficients stœchiométriques des réactifs.
\end{enumerate}
\end{exo}

\begin{exo}[transposer en milieu basique]
On s'intéresse au couple \ce{MnO4^-}/\ce{MnO2} (Mn passe de $+\mathrm{VII}$ à $+\mathrm{IV}$).
\begin{enumerate}[label=\alph*)]
  \item Équilibre la demi-équation en milieu \textbf{acide}.
  \item Transpose-la en milieu \textbf{basique} (avec \ce{OH^-} et \ce{H2O}).
\end{enumerate}
\end{exo}

\begin{exo}[pondérer et comparer les coefficients]
On donne la réaction non pondérée (milieu acide) :
\[
\ce{Sn + NO3^- + H+ -> SnO2 + NO2 + H2O}
\]
(l'étain passe de $0$ à $+\mathrm{IV}$, l'azote de $+\mathrm{V}$ à $+\mathrm{IV}$).
\begin{enumerate}[label=\alph*)]
  \item Écris l'équation pondérée.
  \item La somme des coefficients des réactifs est-elle supérieure, égale ou inférieure à celle des
        produits ?
\end{enumerate}
\end{exo}

\begin{exo}[compter les électrons d'une réaction globale]
La réaction suivante est déjà pondérée :
\[
\ce{2 MnO4^- + 5 H2O2 + 6 H+ -> 2 Mn^{2+} + 5 O2 + 8 H2O}.
\]
\begin{enumerate}[label=\alph*)]
  \item Montre qu'il s'échange \textbf{10} électrons au total.
  \item Quel est l'oxydant ? Quel est le réducteur ?
\end{enumerate}
\end{exo}

\end{document}
